\section{Discussion}
\label{sec:discussion}

\paragraph{Where the speed comes from.}
Single-token decode of a 17\,GB model is bound by memory bandwidth at about 50\,tok/s on this
card, and every engine sits close to that bound without speculation (NInfer 50.5\,tok/s,
llama.cpp 45.9, as reported by the NInfer port). All of the gain therefore comes from three
places: how many tokens a verification pass accepts (the drafter), how cheap the drafter and the
$(k{+}1)$-token verification are (the engine), and how many bytes a pass reads (the weights).
DFlash2 accepts longer runs than the model's own MTP head on structured text, and both
specialized engines make verification cheap. azllm adds the third lever: its weight file is
16\% smaller than the baseline's, and on code it predicts nearly as well, which is unusual for a
smaller file and points to the value of rotation before quantization.

\paragraph{Context is a memory budget.}
On 24\,GB, the weights take 15--18\,GB and everything else competes for the rest: the KV cache,
the drafter and its own cache, the vision tower, prefill buffers and CUDA graphs. Two choices
dominate. 4-bit keys double the context, and a block-diffusion drafter with its own KV cache
costs up to 111K tokens in NInfer. azllm's drafter is small (0.79\,GB), so its fastest
configuration also reaches the full window. Startup checks are not enough: we filled every
retained configuration with real code before calling its context usable.

\paragraph{Precision of a 4-bit KV cache.}
For code, the KV cache matters at long range, where an error in a key shifts attention to the
wrong line. Exact recall of 16-line functions did not separate any format, even at 252K tokens.
Perplexity did: a 4-bit cache costs between 0.16\% (azllm) and 0.82\% (NInfer's $E_8$), against
0.11\% for the baseline's 4-bit values. The difference between the two 4-bit formats is larger
than the difference between 4- and 8-bit keys within azllm, so the format matters more than the
bit count.

\paragraph{Limitations.}
\begin{itemize}
  \item One machine, one GPU, one session. Accuracy runs are single greedy passes; the GSM8K and
        HumanEval differences we report are of the size of their noise.
  \item There is no BF16 reference model: the weight comparison ranks three 4-bit pipelines,
        each with its own arithmetic, not their distance to the original model.
  \item We did not evaluate long agentic sessions with tools. Field reports describe looping
        with four or more MTP drafts in NInfer during long tool-calling runs;
 our NInfer profiles use three, and azllm limits repetition by
        default (\texttt{--repeat-limit 256}). A tool-call smoke test passed on azllm with the
        fixed template.
  \item Speeds were measured with the card at its stock 450\,W limit. The host reset once,
        without a GPU error in the log, during back-to-back compute-bound runs; thereafter a
        guard logged temperatures every 10\,s (peak 80\,$^\circ$C) and never had to pause.
  \item The capacity limits leave each engine's default free memory on the GPU (768\,MiB for
        azllm), but our card also drives the display, whose desktop holds about 1\,GB and grows
        as windows open. Next to it, azllm with \texttt{q8/q4} KV refuses 221{,}184 tokens. Our
        packaged servers keep 2\,GB free and serve 151{,}552 tokens (\texttt{q8/q4}) and
        233{,}472 (\texttt{q4/q4}); the KV cache is reserved whole at start, so the room is for
        the desktop and the CUDA graphs. NInfer and OpenJet lose more, down to 77{,}824 and
        65{,}536 tokens; their limits were measured from the free memory, which the desktop
        moves by up to 400\,MiB between starts.
  \item azllm and NInfer are young projects maintained by few people, and both need weight
        files of their own.
\end{itemize}

\section{Conclusion}
\label{sec:conclusion}

On one RTX~4090, Qwen3.8-27B can be served 2.5$\times$ faster than a tuned llama.cpp setup,
with almost twice its context and vision, at the same precision on code. The engine that
achieves it, azllm, combines a DFlash2 block-diffusion drafter, a specialized CUDA engine and a
rotated 4.3-bit weight format that predicts code within 0.5\% of the baseline's 5.2-bit file.
Its 4-bit KV cache costs 0.16\% in log-likelihood and fits the full 262K-token window with the
vision tower. We release launcher scripts for every configuration measured here, served
behind llama-swap \citep{llamaswap}, which starts the engine a request asks for.
